\section{Cauchy 判据无需预知极限}
\label{sec:12-Real-Analysis-Support/01-Real-Numbers-and-Completeness/02}

要用 \(\varepsilon\)--\(N\) 定义证明 \(x_n\to L\)，你得先知道 \(L\) 是几。写出``对任意 \(\varepsilon>0\)，存在 \(N\)，使 \(n\ge N\) 时 \(|x_n-L|<\varepsilon\)''的前提是你手里已经捏着一个候选极限。

但很多构造场景里，你并不知道极限是谁。你用递推公式生成一项又一项，逐次近似越来越精确，迭代日志里的数字逐渐稳定下来——但你从未见过最终的答案，也没法从递推式里直接读出它。这就像闭着眼睛朝一个目的地走：你能感觉每一步的步长在缩小，能判断自己走得很稳，却说不出目的地的坐标。

你需要一个不依赖``先知道极限''的收敛判据。

\subsection{不跟极限比，跟同伴比}

替代方案出奇的简单：不去比较 \(x_n\) 与某个外部参照物 \(L\)，而是比较数列的项彼此之间。如果数列``想要收敛''，它的尾部各项应该挤在一起——越往后，任意两项之间的距离都应该可以任意小。

形式化：称 \((x_n)\) 为 Cauchy 数列，如果对任意 \(\varepsilon>0\)，存在 \(N\)，使得

\[
m,n\ge N \quad\Longrightarrow\quad |x_m-x_n|<\varepsilon.
\]

这里的量词与收敛定义完全平行：\(\forall\varepsilon\ \exists N\ \forall m,n\ge N\)。差别只在于最后的表达式从 \(|x_n-L|<\varepsilon\) 换成了 \(|x_m-x_n|<\varepsilon\)。极限 \(L\) 从定义中消失了。

这个定义只检查数列的``内聚性''——尾部是否自相靠近，不要求你预先声明它们在靠近什么。

\subsection{收敛一定推出 Cauchy：不依赖完备性}

若 \(x_n\to L\)，给定 \(\varepsilon>0\)，取 \(N\) 使 \(n\ge N\) 时

\[
|x_n-L|<\frac{\varepsilon}{2}.
\]

对任意 \(m,n\ge N\)，三角不等式给出

\[
|x_m-x_n| \le |x_m-L|+|x_n-L| < \frac{\varepsilon}{2}+\frac{\varepsilon}{2} = \varepsilon.
\]

这个方向的证明只有两行，而且不依赖实数完备性。在任何度量空间中——无论是 \(\R\)、\(\R^d\)、还是赋予 sup 距离的函数空间——收敛总是蕴含 Cauchy。收敛强于 Cauchy，或者说 Cauchy 是收敛的必要条件。

真正有意思的是反方向：Cauchy 能否推出收敛？

\subsection{Cauchy 推出收敛需要空间没有空洞}

在 \(\R\) 中，每个 Cauchy 数列都收敛于某个实数。这是实数完备性的另一种等价表述——与上一节的最小上界性质说的是同一件事，只是换了一套语言。上确界版本说``每个有上界的集合有最小上界''；Cauchy 版本说``每个自相靠近的数列有极限''。两种说法从不同角度刻画同一个事实：\(\R\) 没有空洞。

回到有理数的那个裂缝。取有理数截断序列

\[
1,\ 1.4,\ 1.41,\ 1.414,\ 1.4142,\ \ldots
\]

任意两项的差可以任意小——在 \(\Q\) 中它满足 Cauchy 条件。但它在 \(\Q\) 中没有极限，因为 \(\sqrt2\) 不是有理数。有理数世界里的 Cauchy 列追着一个不在场的目标。

在更大的空间 \(\R\) 中看这同一个序列，它收敛到 \(\sqrt2\)。Cauchy 性描述``数列想要收敛''，完备性回答``它想去的地方确实在这个空间里''。不满足完备性的空间，就像一个漏水的容器——有些 Cauchy 列的极限从洞里漏出去了。

\subsection{Cauchy 数列必有界：尾部挤在一起就不可能无限散开}

取 \(\varepsilon=1\)，由 Cauchy 条件存在 \(N\) 使 \(m,n\ge N\) 时 \(|x_m-x_n|<1\)。固定 \(n=N\)，对所有 \(m\ge N\) 有

\[
|x_m| \le |x_N| + 1.
\]

前 \(N-1\) 个项只有有限个，它们的绝对值取最大值与上面的界合并，整个数列就被框在一个有限区间里。

反过来不成立：有界数列未必 Cauchy。\(x_n=(-1)^n\) 有界（始终在 \([-1,1]\) 中），但尾部任意两项可以差 2——取奇数项与偶数项，距离恒为 2，永远压不到 0.1 以下。振荡本身不妨碍 Cauchy，但振幅不衰减的振荡会妨碍。

\subsection{一个常用的判据：相邻差的几何衰减控制整个尾部}

很多时候我们拿不到 \(|x_m-x_n|\) 的直接估计，但可以控制相邻两项的差。若存在常数 \(C>0\) 和 \(q\in(0,1)\) 使

\[
|x_{n+1}-x_n| \le Cq^n,
\]

那么对任意 \(m>n\)，用三角不等式串联起来：

\[
\begin{aligned}
|x_m-x_n|
&\le \sum_{k=n}^{m-1} |x_{k+1}-x_k| \\
&\le C\sum_{k=n}^{\infty} q^k
= \frac{Cq^n}{1-q}.
\end{aligned}
\]

右侧只依赖于 \(n\)，且随 \(n\to\infty\) 趋于零。所以 \((x_n)\) 是 Cauchy 列。在完备空间中，它必收敛。

这个结构出现在三处典型的场景里：

\begin{itemize}
\tightlist
\item 压缩映射迭代：Banach 不动点定理中，\(x_{n+1}=T(x_n)\) 且 \(T\) 是压缩映射，相邻差按压缩因子几何衰减；
\item 数值算法的误差分析：Newton 法在根附近的二次收敛，相邻修正量快速坍缩；
\item 函数级数的 Weierstrass M-test：用几何级数控制逐项上界，从而控制部分和序列的 Cauchy 性。
\end{itemize}

三种场景共享同一个推理骨架：你不需要先知道极限是谁，只要能证明相邻差的尾部可求和，极限的存在就由完备性接管。

\subsection{``相邻差趋于零''不能替代 Cauchy}

一个极其常见的错误是把``相邻两项很接近''当作收敛的证据。许多迭代程序的停止条件是``本次与上次的更新量小于 \(10^{-6}\)''，然后宣称已收敛。

拆开来看，相邻差趋于零——\(|x_{n+1}-x_n|\to 0\)——是 Cauchy 条件的必要但不充分的一部分。调和级数的部分和

\[
x_n = \sum_{k=1}^{n}\frac1k
\]

满足

\[
x_{n+1}-x_n = \frac1{n+1} \to 0,
\]

但 \((x_n)\) 不是 Cauchy 列：对任意 \(n\)，取 \(m=2n\)，有

\[
x_{2n}-x_n = \sum_{k=n+1}^{2n}\frac1k \ge n\cdot\frac1{2n} = \frac12.
\]

无论 \(n\) 多大，总能在后面找到距离不小于 \(1/2\) 的项对。相邻差在缩小，但缩小的速度不够快——每次只加一点点，无穷多次累加起来可以发散。Cauchy 要求的不是相邻项接近，而是\emph{整个尾部挤在一起}，这件事需要累加差被尾部控制和。

调整学习率过小也会造成同样的错觉：参数更新的每一步都很小，但模型可能还在缓慢漂移，只是漂得慢而已。要把``这一步变化很小''解释为``离极限已经很近''，还缺一个链接——要么像几何衰减那样控制整个未来路径的长度，要么有独立的方式把残差与真实误差的关系建立起来。

\subsection{从 Cauchy 列构造实数：完备化}

实数本就可以从有理数的 Cauchy 列等价类构造出来。把两个有理 Cauchy 列视为``代表同一个实数''，如果它们的差趋于零。\(\sqrt2\) 被定义为所有以它为极限的有理 Cauchy 序列构成的等价类——从而在定义层面就把那个空洞填上了。

另一种等价的构造是 Dedekind 分割：把有理数集沿一个``切口''分成左右两部分，切口本身（如果恰好切在无理数上）就定义一个新实数。两种构造都在做同一件事：系统地把 \(\Q\) 的空洞逐一填补，得到没有裂缝的 \(\R\)。

学习基础分析一般不需要展开构造的全部技术细节——那更适合放在集合论或实数模型课程里。但需要知道完备性是一条结构公理，而不是一个需要证明的定理。你在分析里反复使用的收敛定理（单调有界收敛、Cauchy 收敛、Bolzano--Weierstrass、介值定理），都扎根于这条公理。

\subsection{Cauchy 判据怎样延伸到函数空间}

把底层的实数换成函数，度量换成 sup 距离

\[
d_\infty(f,g) = \sup_{x\in E}|f(x)-g(x)|,
\]

函数列是否为 Cauchy 同样有明确含义：尾部所有函数在定义域 \(E\) 的\emph{每一点上同时}互相接近。若函数空间在此距离下完备——即每个一致 Cauchy 列都有空间内的极限函数——就能从逐项近似直接断言极限存在，无需先猜出极限函数的形状。

一致收敛与 Banach 空间的整套语言，正是从这个简单的观察生长出来的：Cauchy 判据让你不必预知极限，完备性保证极限在空间里。
